% ====================================================================
% DATEI: content/15_falsifiability_extended.tex
% Erweiterte Falsifizierbarkeit und experimentelles Programm
% Kern 1.3 — kanonische erweiterte Version
% ====================================================================

\section{Erweiterte Falsifizierbarkeit und experimentelles Programm}
\label{sec:falsifiability-extended}

Dieser Abschnitt stellt den Falsifizierbarkeitsrahmen für die
Ontology of Continua (OC) wieder her und erweitert ihn.
Er formuliert explizite strukturelle Vorhersagen,
Falsifikationskriterien und experimentelle Programme über die Hierarchie
\(K_2\)–\(K_{10}\), ausschließlich auf Basis der Axiome und Definitionen von Core~1.3.3.

Das Ziel ist nicht, ein vollständiges Katalog von Experimenten bereitzustellen, sondern zu zeigen,
dass OC strukturfalsifizierbar ist: Es werden domänenübergreifende Verpflichtungen zu
Schwellenwerten, Zyklen, dimensionalen Übergängen und Randverhalten formuliert,
die prinzipiell durch empirische und theoretische Daten widerlegt werden können.

% ====================================================================
\subsection{Bedeutung der Falsifizierbarkeit für strukturelle Theorien}
% ====================================================================

Für eine strukturelle Theorie ist Falsifizierbarkeit nicht auf einzelne
numerische Vorhersagen beschränkt. Stattdessen tritt Falsifikation ein,
wenn die strukturellen Relationen zwischen
\[
  (\Omega, A, P(t), J(t),
   \Theta, \partial\Omega, C, k(t))
\]
in einer Weise verletzt werden, die mit den Axiomen und Sätzen von OC unvereinbar ist.

\paragraph{Strukturelle vs. phänomenologische Falsifizierbarkeit.}
\begin{itemize}
  \item \emph{Phänomenologische Falsifizierbarkeit} betrifft direkte Abweichungen
        zwischen Modellausgaben und Beobachtungen (z.\,B. falsches Phasendiagramm,
        falscher Skalierungsexponent).
  \item \emph{Strukturelle Falsifizierbarkeit} betrifft Verletzungen universeller
        Restriktionen auf Kontinua: Schwellenwerte, Zyklen, monotone Dimension,
        Einbettungskompatibilität, Kollaps- und Wiederbelebungsbedingungen.
\end{itemize}
OC operiert primär auf der strukturellen Ebene; phänomenologische Modelle müssen
OCs strukturelle Anforderungen erfüllen, um als korrekte Instanziierungen zu gelten.

\paragraph{Kanäle der strukturellen Falsifikation.}
Ein Kontinuumsmodell, das Anspruch auf Instanziierung von OC erhebt, ist falsifizierbar, wenn:
\begin{enumerate}
  \item \textbf{Schwellenwertverletzungen} systematisch auftreten:
        beobachtete Systeme bleiben vital, obwohl sie angebliche
        \(\Theta_{\mathrm{exist}}\), \(\Theta_{\mathrm{stab}}\),
        \(\Theta_{\mathrm{crit}}\) oder \(\Theta_{\mathrm{death}}\) verletzen.
  \item \textbf{Verschwinden von Zyklen} der Vitalität widerspricht:
        Systeme erhalten über lange Zeit eine Organisation ohne irgendein
        identifizierbares Zykluskomplexesystem \(C\), das die Bedingungen von OC erfüllt.
  \item \textbf{Verbotene dimensionsabhängige Übergänge} beobachtet werden:
        Dimension scheint zu steigen, ohne neue Achsen im Einbettungsraum
        zu aktivieren oder ohne Überschreiten eines Dimensionsschwellenwertes
        \(\Theta_{\mathrm{dim}}\).
  \item \textbf{Inkompatible Tupel} notwendig werden:
        erfolgreiche Modelle erfordern für dasselbe System inkompatible Varianten
        des kanonischen OC-Zustandstupels.
\end{enumerate}

\paragraph{Strukturelles vs. empirisches Scheitern.}
Das empirische Scheitern eines konkreten Domänenmodells führt nicht sofort zur Falsifikation von OC;
die strukturelle Theorie ist erst dann in Frage gestellt, wenn der empirische Erfolg
alternativer Modelle Verstöße gegen OCs strukturelle Restriktionen erfordert.
Umgekehrt wird OC direkt falsifiziert, wenn eine Domäne wiederholt Strukturen aufweist,
die OC verbietet (z. B. stabile Kontinua ohne stützende Zyklen).

\paragraph{Cross-domain-Falsifikationslogik.}
Da OC Universalität beansprucht, propagiert eine strukturale Falsifikation in einer Domäne
(wirksam nachgewiesen) auf alle Domänen, die dieselben Axiome verwenden.
Dadurch ist OC \\emph{falsifizierbarer} als eng gefasste Theorien:
\begin{itemize}
  \item Falls die Monotonie der Dimension in einem kontrollierten physischen System versagt,
muss sie auch für kognitive und soziale Kontinua neu bewertet werden;
  \item Falls lebensfähige Protozellen ohne irgendeine Randstruktur nachgewiesen werden,
die mit \(\partial\Omega\) kompatibel ist, ist die gesamte K-Level-Hierarchie betroffen.
\end{itemize}

Der Rest dieses Kapitels formuliert explizite Vorhersagen und Tests.

% ====================================================================
\subsection{Vorhersagen und Tests in der Physik (\texorpdfstring{$K_2$}{K\_2})
% ====================================================================

Physikalische Kontinua bieten den klassischsten Testkontext für OC auf Niveau
\(K_2\).

\paragraph{Strukturelle Vorhersagen.}
\begin{itemize}
  \item \textbf{P2.1 (Kritische Schwellenwerte).}
        Phasenübergänge entsprechen strukturellen Durchtritten von
        \(\Theta_{\mathrm{crit}}\). Kein scharfer Langstrecken-Ordnungsübergang
        kann ohne entsprechende Schwellenfläch in \(\Omega_2\) auftreten.
  \item \textbf{P2.2 (Perkolationsmonotonie).}
        Für konnektivitätsbasierte Systeme mit Steuerparameter \(p\)
erfordert globale Vitalität (\(k_2>0\)) \(p \ge p_c\);
        für \(p<p_c\) kann kein unendlicher Cluster, der zu \(\Omega_2\) kompatibel ist,
        existieren.
  \item \textbf{P2.3 (Monotonie der Dimension).}
        Kein physikalisches Kontinuum kann seine effektive Dimension erhöhen,
        ohne mindestens eine neue Achse im Einbettungsraum \(M_2\)
        zu aktivieren und einen Dimensionsschwellenwert \(\Theta_{\mathrm{dim}}\) zu überschreiten.
  \item \textbf{P2.4 (Randgesteuerter Kollaps).}
        Der Kollaps geordneter Phasen erfolgt über Zerstörung stützender Zyklen
        (z.\,B. Wirbelpaare, Domänenstrukturen) und den Verlust
        eines wohldefinierten \(\partial\Omega_2\), nicht über das spontane
        Verschwinden der Ordnung im Inneren allein.
\end{itemize}

\paragraph{Falsifikationskriterien und Tests.}
\begin{itemize}
  \item \textbf{F2.1.}
        Ein robuster, reproduzierbarer Phasenübergang ohne strukturelle
        Schwellenwerte in der Geometrie des Ordnungsparameters würde P2.1 falsifizieren.
  \item \textbf{F2.2.}
        Die Beobachtung stabiler unendlicher Cluster für \(p<p_c\) in
        gut kontrollierten perkolationsähnlichen Systemen würde P2.2 widersprechen.
  \item \textbf{F2.3.}
        Der Nachweis eines Kontinuums, dessen effektive Dimensionalität steigt
        (z.\,B. von 2D- zu 3D-Verhalten), während alle Achsen von \(M_2\) fixiert bleiben
        und keine neuen Freiheitsgrade aktiviert werden, würde P2.3 sowie den allgemeinen
        Satz zur Monotonie der Dimension falsifizieren.
  \item \textbf{Experimentelles Programm.}
        Präzisionsversuche zur Zerstörung von Kohärenz, Wirbelentfaltung
        und Perkolationsschwellen können genutzt werden, um empirische
        \(\Theta_{\mathrm{crit}}\) zu kartieren und zu prüfen, ob das zugehörige
        \(\partial\Omega_2\) wie von OC vorhergesagt wirkt
        (kritische Verlangsamung, Randlokalisierung des Kollapses).
\end{itemize}

% ====================================================================
\subsection{Vorhersagen und Tests bei Ursprungsfragen des Lebens
            (\texorpdfstring{$K_3$}{K\_3}--\texorpdfstring{$K_4$}{K\_4})}
% ====================================================================

Für Szenarien zum Ursprung des Lebens macht OC explizite Aussagen zur
katalytischen Geschlossenheit, Membranstruktur und Protocell-Vitalität.

\paragraph{Strukturelle Vorhersagen.}
\begin{itemize}
  \item \textbf{P3.1 (Katalytische Geschlossenheit).}
        Langfristig stabile chemische Organisationen auf \(K_3\) benötigen mindestens
        einen RAF-ähnlichen Zykluskomplex; Systeme ohne katalytische Geschlossenheit
        können nicht \(k_3>0\) aufrechterhalten.
  \item \textbf{P3.2 (osmotischer Kollapsschwellenwert).}
        Protocells mit fixer Membranzusammensetzung zeigen eine Radius\-Spannungs\-
        Permeabilitätsbeziehung: Überschreitet die Kombination aus osmotischem
        Gradienten, Krümmung und Permeabilität
        (\(\Theta_{\mathrm{grad}}, \Theta_{\mathrm{perm}}, \Theta_{\mathrm{mem}}\))
        einen kritischen Bereich, ist Kollaps unvermeidlich.
  \item \textbf{P3.3 (Randnotwendigkeit).}
        Protocells, denen eine strukturell identifizierbare Grenze
        \(\partial\Omega_4\) (Membran oder äquivalente Kompartimentstruktur)
        fehlt, können über Zeitfenster, in denen \(k_4>0\) definiert ist,
nicht nichttriviale metabolische Flüsse \(J_{\mathrm{metabolic}}\) aufrechterhalten.
  \item \textbf{P3.4 (Patch- lokalisierter Ausfall).}
        Membran-Kollaps initiiert auf Randbereichen, auf denen lokale Schwellenwerte
        zuerst überschritten werden, nicht gleichmäßig über die gesamte Oberfläche.
\end{itemize}

\paragraph{Falsifikationskriterien und Tests.}
\begin{itemize}
  \item \textbf{F3.1.}
        Der Nachweis chemisch autopoetisch sich erhaltender Organisationen mit
        \(k_3>0\) und ohne identifizierbare katalytische Geschlossenheit oder
        Zykluskomplex würde P3.1 falsifizieren.
  \item \textbf{F3.2.}
        Systematische Verletzung der vorhergesagten osmotischen
        Kollapsbeziehungen (z.\,B. beliebige Gradienten ohne strukturelle Anpassung)
        würde P3.2 widerlegen.
  \item \textbf{F3.3.}
        Beobachtungen langlebiger protocell-ähnlicher Systeme mit robusten
        metabolischen Flüssen, aber ohne mit \(\partial\Omega_4\) kompatible
        Grenzstruktur, würden P3.3 falsifizieren und die Rolle der Grenzen in
        OC abschwächen.
  \item \textbf{Experimentelles Programm.}
        Kontrollierte Experimente zu Protocell-Leckage und -Ruptur in Verbindung mit
        quantitativen Messungen von Gradienten und Membraneigenschaften können den
        strukturellen Status von \(\Theta_{\mathrm{grad}}\),
        \(\Theta_{\mathrm{perm}}\), \(\Theta_{\mathrm{mem}}\) und Patch-Dynamik prüfen.
\end{itemize}

% ====================================================================
\subsection{Vorhersagen und Tests in biologischen Systemen
            (\texorpdfstring{$K_4$}{K\_4}--\texorpdfstring{$K_5$}{K\_5})}
% ====================================================================

Auf biologischer Ebene konzentriert sich OC auf Erregbarkeit, Proto-Spikes und den
Übergang \(K_4 \to K_5\).

\paragraph{Strukturelle Vorhersagen.}
\begin{itemize}
  \item \textbf{P4.1 (Minimale Erregbarkeitsschwelle).}
        Der Übergang von rein metabolischen Protocells zu erregbaren Kontinua
        erfordert eine minimale Erregbarkeitsschwelle \(\Theta_{\mathrm{exc}}\);
        kein stabiles Verhalten auf \(K_5\) kann ohne nichttriviale
        \(\Delta V\)-Achsen und eine wohldefinierte Spike-Schwelle existieren.
  \item \textbf{P4.2 (Proto-Spike-Invarianten).}
        In frühen erregbaren Systemen existieren minimale Invarianten für spike-
        ähnliche Zyklen (z.\,B. minimale Amplitude von \(\Delta V\),
        Erholungszeit \(\tau\)), die den strukturellen Restriktionen von
        \(C_{\mathrm{spike}}\) entsprechen.
  \item \textbf{P4.3 (Erregbarkeit vor komplexer Vererbung).}
        Langfristige Vererbung komplexer Regulationsstrukturen erfordert
        die vorherige Etablierung \(K_5\)-ähnlicher Kanäle und
        Erregbarkeitzyklen; rein \(K_4\)-Systeme können dauerhaft keine
        vererbte regulatorische Komplexität aufrechterhalten,
        ohne auf ein höheres Niveau zu wechseln.
\end{itemize}

\paragraph{Falsifikationskriterien und Tests.}
\begin{itemize}
  \item \textbf{F4.1.}
        Der Nachweis echter \(K_5\)-ähnlicher Signalübertragung (Spike-Züge,
        erregbare Wellen) in Systemen ohne identifizierbare Erregbarkeitachsen
        oder Schwellen würde P4.1 falsifizieren.
  \item \textbf{F4.2.}
        Das Auffinden stabiler, komplexer Vererbungsmechanismen in Protocell-
        Systemen, die strikt auf \(K_4\) beschränkt sind (keine Erregbarkeit,
        keine spike-ähnlichen Zyklen), würde P4.3 herausfordern.
  \item \textbf{Experimentelles Programm.}
        Minimale Ionenkanalsätze, synthetische erregbare Membranen und
        kontrollierte \(\Delta V\)-Instabilitätstests können genutzt werden,
um die strukturelle Landschaft von
        \(\Theta_{\mathrm{exc}}\),\(
        \(\Theta_{\mathrm{refr}}\) und dem resultierenden \(C_{\mathrm{spike}}\) abzutasten.
\end{itemize}

% ====================================================================
\subsection{Vorhersagen und Tests in kognitiven Kontinua
            (\texorpdfstring{$K_6$}{K\_6})}
% ====================================================================

Für kognitive Systeme betont OC Bindungskapazität, Vorhersagefehler und
aus zyklischer Stabilisierung interner Modelle.

\paragraph{Strukturelle Vorhersagen.}
\begin{itemize}
  \item \textbf{P6.1 (Bindungskapazitätsgrenze).}
        Für jedes kognitive Kontinuum gibt es einen endlichen Schwellenwert der
        Bindungskapazität \(\Theta_{\mathrm{bind}}\), oberhalb dessen zusätzliche
        Merkmale nicht mehr in kohärente Repräsentationen integriert werden können,
        ohne dass \(k_6\) kollabiert.
  \item \textbf{P6.2 (Kollaps durch Vorhersagefehler).}
        Es existiert ein Schwellenwert für den Vorhersagefehler
        \(\Theta_{\mathrm{pred}}\), oberhalb dessen kein kohärentes internes
        Modell aufrechterhalten werden kann; kognitiver Kollaps tritt ein,
wenn Vorhersagefehler innerhalb der aktuellen Achsen durch keine Flüsse \(J_6\)
mehr reduziert werden kann.
  \item \textbf{P6.3 (Zyklusbasiertes Kollapsmodell).}
        Kognitiver Kollaps beinhaltet stets die Zerstörung zentraler Zyklen
        (Aufmerksamkeit, Vorhersage, Gedächtnis); es gibt kein Kollaps-Szenario
        mit intaktem \(C_{\max}(K_6)\).
\end{itemize}

\paragraph{Falsifikationskriterien und Tests.}
\begin{itemize}
  \item \textbf{F6.1.}
        Robuste Evidenz für unbeschränkte Bindungskapazität in endlichen Systemen
        (keine beobachtbaren Schwellen, keine Degradierung) würde P6.1 falsifizieren.
  \item \textbf{F6.2.}
        Empirische Regime, in denen Vorhersagefehler divergiert,
doch stabile, kohärente interne Modelle unendlich lange bestehen, widersprechen
        P6.2.
  \item \textbf{F6.3.}
        Der Nachweis eines kognitiven Kollapses, bei dem Aufmerksamkeit-,
        Vorhersage- und Gedächtniszyklen vollständig intakt bleiben,
würde P6.3 und die zyklische Definition des Kollapses widerlegen.
  \item \textbf{Experimentelles Programm.}
        Verhaltens- und neurophysiologische Studien zu Bindungsdegradation
        unter Last, Sättigung von Vorhersagefehlern und Gedächtnisrekonsolidierung
        können die empirischen Gegenstücke von
        \(\Theta_{\mathrm{bind}}\),
        \(\Theta_{\mathrm{pred}}\) und
        \(C_{\mathrm{memory}}\) prüfen.
\end{itemize}

% ====================================================================
\subsection{Vorhersagen und Tests in sozialen/gesellschaftlichen Kontinua
            (\texorpdfstring{$K_7$}{K\_7}--\texorpdfstring{$K_8$}{K\_8})}
% ====================================================================

Auf höheren sozialen Ebenen formuliert OC strukturelle Restriktionen zu Vertrauen,
institutionellen Zyklen und infrastruktureller Stabilität.

\paragraph{Strukturelle Vorhersagen.}
\begin{itemize}
  \item \textbf{P7.1 (Vertrauensschwelle).}
        Stabile soziale Kontinua auf \(K_7\) erfordern Vertrauen oberhalb eines
        Schwellenwertes \(\Theta_{\mathrm{trust}}\); unterhalb dieses Wertes können
        institutionelle Zyklen \(k_7>0\) nicht erhalten bleiben.
  \item \textbf{P7.2 (Verschluss institutioneller Zyklen).}
        Es gibt kein stabiles \(K_7\)-Kontinuum ohne zumindest einen geschlossenen
        institutionellen Zyklus (z.\,B. Steuern–Dienstleistungen–Legitimitäts-Schleife),
der mit OCs Definition von \(C_7\) kompatibel ist.
  \item \textbf{P8.1 (Infrastrukturzyklus-Schwellen).}
        Zivilisationskontinua \(K_8\) benötigen stabile infrastrukturelle Zyklen;
        Kollaps auf \(K_8\) involviert Randversagen in
        \(\partial\Omega_8\) und Perkolation des Versagens über zentrale
        Infrastrukturebenen.
\end{itemize}

\paragraph{Falsifikationskriterien und Tests.}
\begin{itemize}
  \item \textbf{F7.1.}
        Historische oder simulierte Beispiele langfristig stabiler Institutionen
        mit Vertrauen, das dauerhaft unter jedem plausiblen
        \(\Theta_{\mathrm{trust}}\) liegt, würden P7.1 in Frage stellen.
  \item \textbf{F7.2.}
        Der Nachweis dauerhafter, hochkomplexer Gesellschaften ohne irgendeinen
        identifizierbaren geschlossenen Regierungs- oder Ressourcenzyklus würde
        P7.2 falsifizieren.
  \item \textbf{F8.1.}
        Robuste Modelle zivilisatorischen Kollapses, die keine Auflösung
        infrastruktureller Zyklen oder Randversagen in
        \(\partial\Omega_8\) involvieren, würden P8.1 widersprechen.
  \item \textbf{Experimentelles Programm.}
        Quantitative Analysen historischer Daten, Modelle zur Netzwerkanfälligkeit
        und großskalige Simulationen von Vertrauensdynamik und Perkolation der
        Infrastruktur können die strukturelle Rolle von
        \(\Theta_{\mathrm{trust}}\),\(C_{\mathrm{inst}}\) und
        \(\partial\Omega_8\) prüfen.
\end{itemize}

% ====================================================================
\subsection{Theoretische und metatheoretische Vorhersagen
            (\texorpdfstring{$K_9$}{K\_9}--\texorpdfstring{$K_{10}$}{K\_{10}})}
% ====================================================================

OC selbst und andere hochrangige Rahmenwerke können auf theoretischer und
metatheoretischer Ebene getestet werden.

\paragraph{Strukturelle Vorhersagen.}
\begin{itemize}
  \item \textbf{P9.1 (Kohärenzschwellen).}
        Jedes langlebige theoretische Kontinuum \(K_9\) muss Kohärenz- und
        Ausdrucksfähigkeitsschwellen \(\Theta_{\mathrm{coh}}, \Theta_{\mathrm{expr}}\)
erfüllen; anhaltender Erfolg von Theorien, die explizit interne Kohärenz verletzen,
        würde den OC-Rahmen widersprechen.
  \item \textbf{P9.2 (Kaskaden der Inkohärenz).}
        Akkumulierte Inkohärenzen in einem theoretischen Rahmen müssen zu einem
        Kollaps seines Zykluskomplexes führen (Forschungsprogramme, Paradigmen
        Schleifen); anhaltender Fortschritt bei beliebig hoher Inkohärenz ist
        mit OC unvereinbar.
  \item \textbf{P10.1 (Meta-Unvollständigkeit).}
        Jedes metatheoretische Kontinuum \(K_{10}\) ist strukturell unvollständig:
        Mit wachsender Modellmenge muss der strukturelle Spannungswert auf
        \(K_{10}\) schließlich einen \(\Theta_{\mathrm{dim}}\) überschreiten,
        was entweder eine Erweiterung der Achsen oder Kollaps erzwingt.
\end{itemize}

\paragraph{Falsifikationskriterien und Tests.}
\begin{itemize}
  \item \textbf{F9.1.}
        Ein vollständig intern konsistentes, abgeschlossenes Theoriemodell,
das alle relevanten Modelle erfasst, ohne je Expansion oder Revision zu verlangen,
um dennoch unbegrenzten empirischen Fortschritt zu stützen, würde
        P9.2 und P10.1 herausfordern.
  \item \textbf{F9.2.}
        Der Nachweis von theoretischen Rahmenwerken, die nach OC-Standards
        inkohärent bleiben, aber dennoch stabile, kumulative,
domänenübergreifende prädiktive Erfolge zeigen, würde die strukturelle Verbindung
        zwischen Kohärenz und \(k_9\) falsifizieren.
  \item \textbf{Experimentelles Programm.}
        Formale Analysen von Modellfamilien, logische Konsistenzprüfungen und
        metatheoretische Untersuchungen des Theorienwandels können zur Prüfung der
        Existenz und des Verhaltens von
        \(\Theta_{\mathrm{expr}}\), \(\Theta_{\mathrm{coh}}\) und zugehörigen
        Zyklen auf \(K_9\)–\(K_{10}\) genutzt werden.
\end{itemize}

% ====================================================================
\subsection{Meta-Kriterien und domänenübergreifende Validierung}
% ====================================================================

Schließlich formuliert OC metathematische Kriterien, die Falsifizierbarkeit über
Domänen und K-Ebenen koppeln.

\paragraph{Prädiktive Invarianten über K-Ebenen hinweg.}
Mehrere strukturelle Relationen werden als universell vorhergesagt:
\begin{itemize}
  \item Existenz mindestens eines stützenden Zyklus \(C_{\mathrm{support}}\)
        für jedes aktive Kontinuum;
  \item Monotonie der Dimension unter Geburtsoperatoren;
  \item Notwendigkeit der Kompatibilität des Einbettungsraums
        \(A(K_x)\subseteq A(M_x)\) und der Schwellenwertkompatibilität mit \(M_x\);
  \item Kollapssignaturen: kritische Verlangsamung, Randlokalisierung,
        Zykluszerfall, Divergenz der strukturellen Spannung.
\end{itemize}
Das Auftreten von Verletzungen dieser Invarianten auf irgendeiner gut
verstandenenen K-Ebene würde sowohl aufwärts als auch abwärts in der Hierarchie
propagieren.

\paragraph{Algorithmische Strukturtests.}
Core~1.2 und nachfolgende Arbeiten sollen algorithmische Tests entwickeln für:
\begin{itemize}
  \item \(\Omega\)-Konsistenz (nichtleerer zulässiger Bereich gegeben Schwellenwerte),
  \item Zyklenvitalität (Existenz von Zykluskomplexen \(C\) bei gegebenen
        Flüssen),
  \item Einbettungskompatibilität (Achsen und Schwellenwerte vs. \(M_x\)).
\end{itemize}
Diese Tests können auf konkrete Modelle in Physik, Chemie,
Biologie, Kognition und Sozialwissenschaften angewendet werden, um zu prüfen,
ob sie legitime OC-Instanziierungen sind.

\paragraph{Cross-domain-Falsifikationslogik.}
Weil Axiome auf den K-Ebenen geteilt werden, besitzt Falsifikation eine
\emph{cross-domain}-Struktur:
\begin{itemize}
  \item Fällt ein strukturelles Axiom in einem streng kontrollierten physischen
        Kontext aus, müssen alle daraus abgeleiteten höherstufigen Vorhersagen
        neu bewertet werden;
  \item umgekehrt erhöht eine konsistente Bestätigung desselben strukturellen
        Motivs (z.\,B. randgesteuerter Kollaps) in mehreren Domänen das Vertrauen
        in das zugehörige OC-Axiom.
\end{itemize}

\paragraph{Integrierte Validierungspipeline.}
Core~1.3.3 liefert die strukturelle Grundlage für eine explizite
Validierungspipeline in Core~1.2:
\begin{enumerate}
  \item Domänenmodelle auswählen und ihre Kontinuumskomponenten identifizieren
        \((\Omega, A, P, J, \Theta, \partial\Omega, C, k)\);
  \item lokale strukturelle Vorhersagen testen (Schwellen, Zyklen, Grenzen);
  \item invarianten über Ebenen hinweg testen (dimensinale Monotonie,
        Kompatibilität von Einbettungen);
  \item Abbildung zwischen Domänentheorien und K-Ebenen aktualisieren.
\end{enumerate}

OC ist daher kein rein konzeptionelles Framework: Es ist darauf ausgelegt,
empirisch beschränkt und~— im Prinzip~— strukturfalsifizierbar über die gesamte
Hierarchie \(K_0\)–\(K_{10}\).

% END OF FILE
